% Options for packages loaded elsewhere
\PassOptionsToPackage{unicode}{hyperref}
\PassOptionsToPackage{hyphens}{url}


\documentclass[]{article}

\usepackage[T1]{fontenc}
\usepackage{lmodern}
\usepackage{graphicx}

\usepackage[top=2.54cm, bottom=2.54cm]{geometry}

\usepackage{setspace}
\onehalfspacing

\usepackage{indentfirst}
\setlength{\parindent}{1.27cm}
\setlength{\parskip}{0pt}


\usepackage{titlesec}
\titleformat{\section}[block]{\normalfont\bfseries\filcenter}{}{0em}{}
\setcounter{secnumdepth}{-\maxdimen} 

\usepackage{hanging}
\usepackage{float}

\usepackage{xurl} 
\usepackage[hidelinks]{hyperref}
\urlstyle{same}


\newcommand{\imagen}[1]{
    \begin{figure}[H]
        \centering
        \includegraphics[width=\textwidth]{#1}
    \end{figure}
}


\begin{document}
    \begin{center}
        \vspace*{2cm}
        \textbf{Tarea Corta 05: Capítulos de \emph{Little Book of Semaphores}}
        \url{https://github.com/sergiopb06/Tareas_Paralela_2026.git}
        
        Sergio Peralta Barboza

        Escuela de Ciencias de la Computación e Informática, Universidad de Costa Rica

        CI-0117 Programación Paralela y Concurrente, Grupo 01

        Ing. Sleyter Angulo, M.Sc

        Jueves 8 de octubre de 2026
    \end{center}
\newpage


\section*{Readers-writers problem}

    El readers-writers problem ocurre cuando una estructura de datos
    compartida es leída y modificada por hilos concurrentes. Mientras se
    escribe, hay que impedir que otros hilos puedan acceder a la estructura,
    ya que podrían leer datos inconsistentes o erróneos. El problema indica
    que cualquier cantidad de lectores puede estar en la sección crítica del
    programa a la vez. Mientras los escritores deben tener acceso exclusivo
    a la sección crítica, es decir, solo puede haber un escritor en la
    estructura, nadie más. Esto se llama exclusión mutua categórica; la
    presencia de una categoría (lectores) excluye a otra (escritores), pero
    no a sus propios miembros.
    
    Para la solución a este problema se propone lo siguiente:
    
    \imagen{./12619979f0363bbc10758ba58951b3c2c3ece808.png}
    Este es el código para los escritores; significa que si la sala está
    vacía, el escritor puede entrar y excluir a todos los otros hilos de
    entrar; al salir, se asegura de que la sala está vacía. Para los
    lectores:
    
    \imagen{./1a548f8e7e14257ebd38eff99ee74a14fda865fa.png}
    Este código funciona de la siguiente forma: cuando un lector quiere
    entrar, se verifica si es el primero; si lo es, señala que el cuarto no
    está vacío. Sin embargo, igualmente permite que otros lectores sigan
    entrando. Conforme salen, si un lector es el último en salir, señala que
    el cuarto ahora está vacío. Si el lector intenta entrar mientras hay un
    escritor en el cuarto, espera en el semáforo del cuarto mientras tiene
    el mutex. Esto significa que todos los otros lectores que quieran entrar
    van a esperar en el mutex, no en el semáforo. El capítulo explica que
    este patrón es común y que se puede empaquetar en una clase
    \emph{Lightswitch:} el primero que entra prende las luces, el último que
    sale las apaga:
    
    \imagen{./57c19d305fd96d866700d7762bca014a3bbc4ffd.png}
    Con esta clase, el código queda más simple:
    
    \imagen{./79c8367efd751136c359f528bcb96c12243f8286.png}
    
    Esta solución presenta un problema, \emph{starvation}: si un escritor
    llega mientras hay lectores en cola, mientras sigan llegando lectores
    antes que el último se vaya, el escritor nunca logra entrar a la cola,
    haciendo que espere indefinidamente para escribir. Por eso se presenta
    el problema de \emph{No-starve}, para modificar el código y que el
    escritor no sufra de \emph{starvation}.


\section*{No-starve readers-writers}

    Para la solución de este problema se utiliza una barra de acceso
    (torniquete). Cuando un lector llega, si el torniquete está abierto,
    sigue con normalidad y entra a la sección crítica. Si un escritor llega
    en la cola, cierra el torniquete y espera a que el cuarto esté vacío. El
    escritor espera a que todos los lectores que estén dentro del cuarto
    terminen su trabajo, y cuando lo hacen, entra al cuarto. Los nuevos
    lectores que llegan esperan su turno hasta que se abra el torniquete, el
    cual es abierto por el escritor una vez termine su trabajo y salga del
    cuarto. Esto asegura que el escritor avance su turno y que los lectores
    esperen mientras lo hace. Los lectores no tocan el torniquete, solo
    pasan por él. La solución propuesta es la siguiente:
    
    \imagen{./a262dc38ab13f6cdcc6cf855da083e4bf21d5a41.png}
    \imagen{./d820d893b596b7083a1936dae934052c1425ae1e.png}
    
    Sin embargo, esta solución solo resuelve que al menos un escritor avance
    en la cola. Pero no garantiza que todos los escritores avancen; cuando
    el escritor termina, es decisión del planificador quién despierta; un
    lector puede colarse delante de un escritor. Por eso hay que crear una
    solución que le dé prioridad a los escritores.

\section*{Writer-priority readers-writers}

    Para la solución de este problema, se reemplaza el semáforo de la
    habitación y el torniquete por un semáforo que indica cuándo hay
    lectores y otro de cuándo hay escritores. El lector pasa por
    \emph{noReaders}, igual que en el torniquete, y mientras tiene
    \emph{noReaders}, si es el primero, agarra \emph{noWriters}. Esto
    garantiza que, mientras haya lectores dentro de la habitación, no pueden
    entrar escritores. El escritor primero agarra \emph{noReaders}; esto
    hace que, cuando llega un escritor, bloquee que entren los lectores
    hasta que todos los escritores terminen su trabajo. Los demás escritores
    que lleguen no tocan el semáforo de los lectores, van esperando su turno
    para hacer su trabajo y, cuando ya no hay escritores en la cola, liberan
    \emph{noReaders} y permiten a los lectores avanzar. La solución
    propuesta es:
    
    \imagen{./390928aba978bc85f7165197da3c6ef2f044a487.png}
    \imagen{./d269e73feee2cdc30a10f07d255bffe6bc04802b.png}
    \imagen{./f1750fc7d3508043fee08368e66868b8e6241762.png}
    
    Esta solución crea el problema inverso; ahora los lectores pueden sufrir de
    \emph{starvation}. La implementación que se use depende de la
    aplicación.
    
\section*{No-starve mutex}

    En la sección anterior, el problema era \emph{categorical starvation}.
    Una categoría (lectores) dejaba sin turno a otra (escritores). Ahora el
    problema es \emph{thread starvation}, que un hilo espere indefinidamente
    mientras otros avanzan. Para resolverlo se exige \emph{bounded waiting};
    el tiempo que un hilo espera en un semáforo debe ser demostrablemente
    finito. Para demostrar que no existe \emph{starvation}, se realizan las
    siguientes suposiciones sobre quién decide el orden del hilo:
    
    \begin{enumerate}
    \def\labelenumi{\arabic{enumi}.}
    \item
      Si solo hay un hilo listo para correr, el planificador debe dejarlo
      correr.
    \item
      Si un hilo está listo para correr, el tiempo que espera para hacerlo
      está acotado.
    \item
      Si hay hilos esperando en un semáforo y alguien hace signal, uno de
      los que esperan debe despertar.
    \item
      Si un hilo espera en un semáforo, el número de hilos que despertarán
      antes que él está acotado.
    \end{enumerate}
    
    Los semáforos que tienen la propiedad cuatro se llaman semáforo fuerte
    (\emph{strong}) y los que tienen solo la propiedad tres se llaman
    semáforo débil (\emph{weak}). Para resolver una conjetura de Dijkstra,
    que propuso que no se puede resolver la exclusión mutua sin
    \emph{starvation} usando solo semáforos débiles. En 1979, J.M. Morris lo
    refutó suponiendo un número finito de hilos.
    
    Su algoritmo usa dos torniquetes (t1 y t2) que forman dos salas de
    espera antes de la sección crítica; funciona en dos fases. En la
    primera, t1 está abierto y t2 cerrado, así que los hilos se acumulan en
    la segunda sala. En la segunda, t1 se cierra y t2 se abre; los hilos
    acumulados pasan a la sección crítica. Los contadores room1 y room2
    indican cuántos hilos hay en cada sala y un mutex protege room1. El
    último hilo en cruzar la sala 1 abre t2; el último en salir de la sala 2
    lo vuelve a abrir. La solución descrita es la siguiente:
    
    \imagen{./4d2685eeec27d4a91ba15d23f71948ba7264098f.png}

\section*{Cigarrete Smokers problem}

    El siguiente caso, \emph{Cigarette Smokers Problem}, fue presentado por
    Suhas Patil. El problema dice que hay cuatro hilos: un agente y tres
    fumadores. Para fumar, un fumador necesita tres ingredientes: tabaco,
    papel y fósforos. Se asume que:
    
    \begin{itemize}
    \item
      Cada fumador tiene una provisión infinita de un solo ingrediente (uno
      tiene tabaco, otro papel, otro fosforos).
    \item
      El agente tiene provisiones infinitas de los tres. Repetidamente,
      elige dos ingredientes distintos al azar y los pone en la mesa.
    \item
      El fumador que tiene el ingrediente que falta debe tomar los dos de la
      mesa, hacer el cigarro, fumar y avisarle al agente para que ponga
      otros dos.
    \end{itemize}
    
    El agente representa un sistema operativo que asigna recursos; los
    fumadores, aplicaciones que necesitan dichos recursos. Existen tres
    versiones del problema, pero se resolverá la versión interesante, que
    dice que no se puede modificar el código del agente, el cual es:
    
    \imagen{./44b78f08607f86493fd3635819c659e2ccbe5891.png}
    
    El agente son tres hilos (A, B, C) que esperan en el semáforo de
    \emph{agentSem}. Cada vez que se hace \emph{signal} de este, uno de los
    tres despierta y pone dos ingredientes en la mesa utilizando los
    semáforos de ingredientes.
    
    \imagen{./2f432fad0800aa3ff84c77ebd44785e51032e1d6.png}
    
    Sin embargo, la solución natural no funciona, la cual se basa en que
    cada fumador espere sus dos ingredientes. Falla por un \emph{deadlock}.
    Por ejemplo, si el agente pone tabaco y papel, el fumador con fósforos
    espera tabaco y papel; puede avanzar. Pero el fumador con tabaco espera
    papel; también puede despertar. Un fumador toma tabaco y otro papel.
    Entonces, el primero queda esperando el papel y el segundo queda
    esperando fósforos, que nunca van a llegar; se crea un \emph{deadlock}.
    
    La solución que propuso Parnas consiste en agregar tres hilos
    auxiliares, pusher, cada uno asociado a un ingrediente. Cuando el agente
    pone un ingrediente, el pusher correspondiente revisa con un
    \emph{mutex} si el otro ingrediente que completa el par ya está en la
    mesa. Si es así, señala al fumador adecuado y limpia la variable
    booleana; si no, solo registra que su ingrediente está disponible. Los
    fumadores ahora solo esperan por un único semáforo, que solo se activa
    cuando los dos ingredientes que necesita están listos:
    
    \imagen{./6ff42e8d9625cc65557b3a2102311ee4bf503004.png}
    \imagen{./bf41038b2657956f53933597575149bc19b7f0d6.png}
    \imagen{./10be9ea54de8a8afc9f67a6ce7873c5feff141ba.png}
    
\section*{Generalized Smokers problem}
    
    En la versión anterior, el agente ponía dos ingredientes en la mesa y
    esperaba a que un fumador terminara antes de poner otros. Eso
    garantizaba que en la mesa nunca hubiera más de un par. Parnas propone
    quitar esa espera: el agente sigue poniendo ingredientes sin esperar a
    nadie. Ahora puede haber repeticiones de los ingredientes que hay en la
    mesa. Como el agente ya no espera, \emph{agentSem} ya no se usa. Además,
    se dejan de usar valores booleanos para los ingredientes; ahora se
    cuentan las unidades.
    
    \imagen{./0ba91c05a46a74dc856b857afdb63ea9a381b32c.png}
    
    El resto de la estructura se mantiene: cada \emph{pusher} espera la señal
    del agente, toma el \emph{mutex} y revisa el estado. Si con sus
    ingredientes y los de la mesa se completa la pareja, decrementa el
    contador del ingrediente y despierta al fumador correspondiente.
    
    Cada vez que un agente corre, es como si creara dos \emph{pushers} con
    un ingrediente cada uno y los enviara a una habitación con los demás.
    Con el \emph{mutex,} entran uno a la vez a una habitación donde hay una
    mesa y tres fumadores dormidos. Cada \emph{pusher} revisa la mesa y, si
    puede completar la pareja, se lo lleva y despierta al fumador correcto;
    si no, deja su ingrediente y se va sin despertar a nadie. Se menciona
    que este es un patrón que aparece en otros problemas, que llaman
    \emph{scoreboard}. Las variables funcionan para mantener un control del
    estado del sistema, mientras cada hilo pasa por el \emph{mutex}, revisa
    el estado (\emph{scoreboard}) y reacciona en consecuencia, sin necesidad
    de comunicarse con otros hilos. El código modificado para el
    \emph{pusher A} es:
    
    \imagen{./56ccc3d2c959c4742fc7288eda65ef28821a72dc.png}

\section*{Barbershop problem}

    Este problema fue propuesto por Dijkstra. El problema de la barbería
    plantea un local con una sala de espera de n sillas y un cuarto con la
    silla del barbero. Si no hay clientes, el barbero duerme; si un cliente
    llega y todas las sillas están ocupadas, se va; si el barbero está
    ocupado, pero hay sillas libres, el cliente espera sentado; y si el
    barbero duerme, el cliente lo despierta. Se dice que los clientes
    invocan \emph{getHairCut}, que un cliente que llega a una barbería llena
    invoca \emph{balk} (se va), y que el barbero invoca \emph{cutHair}.
    Además, que mientras el barbero invoca \emph{cutHair}, exactamente un
    hilo debe estar invocando \emph{getHairCut}.
    
    Se usa un contador customers con el número de clientes en la barbería,
    protegido por un \emph{mutex}, y un límite de n, que son las sillas de
    espera más la silla del barbero. Además, hay cuatro semáforos que forman
    dos \emph{rendezvous}: customer y barber para comenzar el corte, y
    \emph{customerDone} y \emph{barberDone} para terminarlo. Los nombres
    siguen la convención de que \emph{customer.wait()} significa ``esperar a
    un cliente''.
    
    \imagen{./fe80adb7c471fbfc39c1db48c5dbad1de07cd068.png}
    
    La solución del cliente es que consulta el marcador (\emph{scoreboard})
    bajo el \emph{mutex}; si la barbería está llena, libera el \emph{mutex}
    y se va; si no, incrementa el contador. Luego avisa al barbero y espero
    a que lo llame. Después recibe el corte y avisa que terminó, espera la
    confirmación del barbero y decrementa el contador.
    
    \imagen{./0816dc7b8bdc5274245b2e04f4211a4180300708.png}
    
    La solución del barbero es esperar a un cliente con
    \emph{customer.wait()}, lo llama con \emph{barber.signal()}, le corta el
    pelo, espera a que el cliente termine con \emph{customerDone.wait()} y
    confirma con \emph{barberDone.signal()}. Si hay clientes esperando, la
    siguiente iteración no duerme.
    
    \imagen{./e6e11a5ed1e78546637b7c9560128a31c8f1c788.png}
    
\section*{FIFO Barbershop}

    La solución anterior de la barbería no garantiza que los clientes sean
    atendidos en el orden en que llegan. Hasta n clientes pueden pasar el
    torniquete, avisar al barbero y quedarse esperando en el mismo semáforo.
    Cuando el barbero señala, cualquiera de los que espera puede despertar.
    El problema es modificar la solución para que los clientes sean
    atendidos en el orden en que llegan.\\
    La solución consiste en que, en lugar de que todos los clientes esperen
    un semáforo compartido, cada cliente tiene su propio semáforo. Al pasar
    el torniquete, el cliente deja ese semáforo en una cola. Cuando el
    barbero está listo, saca de la cola el semáforo del primer cliente. De
    este modo, el barbero es el que decide a quién atender y no el sistema.
    
    \imagen{./b14a35ea5ecafa7067bc653cd8fdb585b74df935.png}
    
    El código del cliente crea su semáforo, toma el mutex, verifica la
    capacidad y, si hay lugar, incrementa el contador y agrega su semáforo a
    la cola. Después avisa al barbero y espera en su propio semáforo. Cuando
    es atendido, el resto queda igual.
    
    \imagen{./d6fd91101b5fa7373aeb22bfdb9dfd5a4e251bcd.png}
    
    El código del barbero espera a un cliente con customer.wait(), toma el
    mutex, saca el primer semáforo de la cola, suelta el mutex y lo señala.
    El resto queda igual.
    
    \imagen{./21c51460d53a925216bb67e3bbd4147df24eed83.png}

\section*{Hilzer's Barbershop problem}

    Este problema presentado por William Stallings es una versión más
    complicada del problema de la barbería. La barbería tiene tres sillas,
    tres barberos y una sala de espera con un sillón para cuatro personas
    más espacio para esperar parado. La capacidad máxima de clientes
    esperando en la tienda es de 20, y quienes lleguen cuando está llena no
    entran. Los clientes esperan de pie, luego se sientan en el sillón y
    luego pasan a la silla para el corte de pelo. El que más tiempo lleva en
    cada etapa es el primero en avanzar, al igual que en \emph{FIFO}. Al
    terminar, cualquier barbero cobra, pero solo hay una caja, entonces solo
    se cobra a un cliente a la vez.
    
    Existen las siguientes restricciones:
    
    \begin{itemize}
    \item
      Los clientes invocan estas funciones en orden: \emph{enterShop,
      sitOnSofa, getHairCut, pay.}
    \item
      Los barberos invocan: \emph{cutHair} y \emph{acceptPayment.}
    \item
      Los clientes no pueden invocar \emph{enterShop} si la tienda está
      llena.
    \item
      Si el sillón está lleno, un cliente que llega no puede invocar
      \emph{sitOnSofa}.
    \item
      Cuando un cliente invoca \emph{getHairCut}, debe haber un barbero
      concurrente ejecutando \emph{cutHair}, y viceversa.
    \item
      Debe ser posible que hasta tres clientes concurrentes ejecuten
      \emph{getHairCut} y que hasta tres barberos ejecuten \emph{cutHair}.
    \item
      Los clientes deben pagar (\emph{pay}) antes de que un barbero pueda
      \emph{acceptPayment}.
    \item
      El barbero debe \emph{acceptPayment} antes de que el cliente pueda
      irse.
    \end{itemize}
    
    La solución aplica dos veces la técnica \emph{FIFO}: cada cliente tiene
    un semáforo propio por cada cola, de modo que el barbero decide a quién
    llamar y se respeta el orden. En la \emph{queue1}, el barbero saca al
    cliente más antiguo y espera a que el siguiente se siente en el sillón.
    En el \emph{queue2}, el cliente que lleva más tiempo en el sillón lo
    manda a la silla. Cada barbero atiende un cliente por vuelta, de modo
    que hay tres cortes concurrentes. El cliente verifica la capacidad y se
    mete en \emph{queue1}, avisa y espera su turno, se sienta en el sillón,
    avisa al barbero, se anota en \emph{queue2}, espera la silla, libera su
    lugar en el sillón, se corta el pelo y paga. El barbero espera a un
    cliente de la \emph{queue1}, lo deja pasar al sillón, luego espera a uno
    de la \emph{queue2}, lo llama a la silla, le corta el pelo y acepta y
    confirma el pago.
    
    \imagen{./527895935f52ec64399d2130f78cd378334469c8.png}
    \imagen{./a2abf6e5e5d873082fd23d447ecc63d87bb58f6b.png}
    \imagen{./e184ad98327fea4b9177d504ec86ddca15939b9a.png}
    
\end{document}
